\section{线性差分方程组与状态变量分析}\label{sec:Difference_Equation}

\subsection{线性差分方程组与常系数线性差分方程组}\label{sub:linear_difference_equation}

\begin{definition}[一阶线性差分方程组]
    设$A[n]$为$N\times N$的实矩阵值函数，$\mathbf{f}[n]$为$N$维实向量值函数，则形如
    \begin{align*}
        \mathbf{y}[n+1] = A[n] \mathbf{y}[n] + \mathbf{f}[n]
    \end{align*}
    的方程组称为\textbf{一阶线性（前向）差分方程组}，其中$\mathbf{f}[n]$为\textbf{非齐次项}。

    特别地，如果$\mathbf{f}[n]=\mathbf{0}$，则称之为\textbf{齐次线性差分方程组}；如果$A[n]$为常数矩阵，则称为\textbf{常系数线性差分方程组}。

    将方程的个数$N$称为方程组的\textbf{阶数}。
\end{definition}

不难看出，类似于\ref{sub:compare_to_high_order}中微分方程组与微分方程的关系，\textbf{差分方程组也能与同阶的（多个）差分方程相互转化}。考虑前向差分方程组，是因为方程组不涉及因果性的问题，而前向差分方程组是更简单、常见的。

任给正整数$k$，可以通过迭代法求出\begin{align*}
    \mathbf{y}[k] &= A[k-1] \mathbf{y}[k-1] + \mathbf{f}[k-1]\\
    & = A[k-1] (A[k-2] \mathbf{y}[k-2] + \mathbf{f}[k-2]) + \mathbf{f}[k-1]\\
    & = A[k-1] A[k-2] \mathbf{y}[k-2] + A[k-1] \mathbf{f}[k-2] + \mathbf{f}[k-1]\\
    & \cdots\\
    & = A[k-1] A[k-2] \cdots A[0] \mathbf{y}[0] + \sum_{i=0}^{k-1} A[k-1] A[k-2] \cdots A[i+1] \mathbf{f}[i]
\end{align*}
因此，给定初始值$\mathbf{y}[0]$，就可以\textbf{唯一确定}任意时刻的$\mathbf{y}[k]$。特别地，如果$A[n]$为常数矩阵，则\begin{align*}
    \mathbf{y}[k] = A^k \mathbf{y}[0] + \sum_{i=0}^{k-1} A^{k-1-i} \mathbf{f}[i]
\end{align*}

利用解的存在性与唯一性，可以仿照\ref{sub:solution_of_linear_equation_system}中的方法进一步说明：
\begin{theorem}
    $N$阶线性差分方程组的解空间$S$构成$N$维实线性空间。
\end{theorem}
\begin{proof}
    线性性是显然的。我们构造一个同构映射来说明解空间的维数是$N$。

    定义$\mathbf{H}:\mathbb{R}^N\to S,\mathbf{y}_0\mapsto \mathbf{y}[n]$，其中$\mathbf{y}[n]$是满足初始条件$\mathbf{y}[0]=\mathbf{y}_0$的解。下验证$\mathbf{H}$是一个线性空间同构映射：

    首先，$\mathbf{H}$是线性的。$\forall \mathbf{y}_0,\mathbf{z}_0\in\mathbb{R}^N,\alpha,\beta\in\mathbb{R}$，$\alpha \mathbf{H}(\mathbf{y}_0) + \beta \mathbf{H}(\mathbf{z}_0)$是满足初始条件$\alpha \mathbf{y}_0 + \beta \mathbf{z}_0$的解，根据解的唯一性，
    $$\mathbf{H}(\alpha \mathbf{y}_0 + \beta \mathbf{z}_0) = \alpha \mathbf{H}(\mathbf{y}_0) + \beta \mathbf{H}(\mathbf{z}_0)$$

    其次，$\mathbf{H}$是单射。设$\mathbf{y}_0\neq\mathbf{z}_0$，它们对应的解分别为$\mathbf{y}=\mathbf{H}(\mathbf{y}_0)$和$\mathbf{z}=\mathbf{H}(\mathbf{z}_0)$，则$\mathbf{y}[0]\neq\mathbf{z}[0]$，因此$\mathbf{y}\neq\mathbf{z}$。

    最后，$\mathbf{H}$是满射。对于任意$\mathbf{y}[n]\in S$，$\mathbf{y}_0=\mathbf{y}[0]$满足$\mathbf{H}(\mathbf{y}_0)=\mathbf{y}[n]$。

    综上，$\mathbf{H}$是一个线性空间的同构映射，解空间的维数为$\dim \mathbb{R}^N=N$。
\end{proof}

由于高阶线性差分方程总能转化为同阶的线性差分方程组，因此高阶线性差分方程的解空间也是一个维数为$N$的线性空间，这就解释了为什么上一节中我们总是给出$N$个初始值来唯一确定一个高阶线性差分方程的解。

前面说明解的存在性与唯一性时，已经给出了通解的表达式，只是在$A[n]$不为常数值矩阵时，没有什么比较好的办法能够简化解的计算，除非$A[n]$具有某些特殊的结构，例如对角矩阵、上三角矩阵等，但这个问题在常系数线性差分方程组中是易于解决的。

对于常系数线性差分方程组，我们也已经给出了通解的表达式：
$$\mathbf{y}[k] = A^k \mathbf{y}[0] + \sum_{i=0}^{k-1} A^{k-1-i} \mathbf{f}[i]$$
因此，求解常系数线性差分方程组的关键在于求出$A^k$的表达式。在多数问题中，我们只需要求出有限的几个$\mathbf{y}[k]$，除了直接迭代，通常有两种快速算法：Jordan分解算法和Cayley-Hamilton定理递推算法。

假设已知矩阵$A$的Jordan分解$A=PJP^{-1}$，其中$J$是一个块对角矩阵，就可以很容易地求出$J^k$的表达式，从而
$$A^k =(PJP^{-1})^k = P J^k P^{-1}$$

Cayley-Hamilton定理（见\ref{sub:eigenvector}）的内容是，设矩阵$A$的特征多项式为
$$p(\lambda)=\det(\lambda I - A)=\sum_{i=0}^{N} c_i \lambda^i$$
则
$$p(A)=\sum_{i=0}^{N} c_i A^i=0$$
因此，$A^N$可以表示为$A^0,A^1,\cdots,A^{N-1}$的线性组合，即
$$A^N = -\sum_{i=0}^{N-1} c_i A^i$$
这样，对于任意的$k\geq N$，我们都可以将$A^k$表示为$A^0=I,A^1,\cdots,A^{N-1}$的线性组合，并递归地求出线性组合的系数，从而求出$\mathbf{y}[k]$的表达式。

下面给出并比较三种算法的\textbf{时间复杂度}。

(1)直接迭代

直接利用递推式进行计算：
$$\mathbf{y}[k+1] = A \mathbf{y}[k] + \mathbf{f}[k]$$
每一步都需要计算复杂度为$O(N^2)$的矩阵乘法，因此递归地求取$\mathbf{y}[k]$的复杂度为$O(k N^2)$。

这一方法没有预处理的时间成本，并且同样适用于$A[n]$为非恒定矩阵的情况，但当$k$较大时，计算效率较低。

(2)Jordan分解算法

在这种算法中，我们首先需要将$A$做Jordan分解，得到$A=PJP^{-1}$，这个过程包括求特征值和特征向量，以及求矩阵$P$的逆，复杂度为$O(N^3)$。然后，我们需要计算$J^k$，由于$J$是一个块对角矩阵，$J^k$的每个块都有已知的解析表达式，因此$J^k$的计算复杂度为$O(N^2)$。因此，Jordan分解算法的预处理复杂度为$O(N^3)$。

对于齐次方程，每求一个$\mathbf{y}[k]$，都要计算矩阵乘法$A\mathbf{y}[k] = P J^k P^{-1} \mathbf{y}[0]$，复杂度为$O(N^3)$，因此每求一个$\mathbf{y}[k]$的复杂度为$O(N^3)$。

对于非齐次方程，由于需要计算$\sum_{i=0}^{k-1} A^{k-1-i} \mathbf{f}[i]$，每次矩阵乘法$A^{k-1-i} \mathbf{f}[i]$的复杂度为$O(N^2)$，因此每求一个$\mathbf{y}[k]$的复杂度为$O(N^3)+O(k N^2)$。

(3)Cayley-Hamilton定理递推算法

为了计算矩阵的特征多项式，通常使用Hessenberg分解算法，复杂度为$O(N^3)$；通过矩阵乘法计算$A^i,i=2,3,\cdots,N-1$，复杂度为$O(N^3)$。在预处理时，还需计算$\mathbf{y}[0],\mathbf{y}[1],\cdots,\mathbf{y}[N-1]$，以免后面将$A^k$表示为$A^0=I,A^1,\cdots,A^{N-1}$的线性组合之后进行重复的矩阵乘法运算，这个过程的复杂度为$O(N^3)$。因此这种算法的预处理复杂度为$O(N^3)$。

对于齐次方程，$\mathbf{y}[k]$可以表示为$\mathbf{y}[0],\mathbf{y}[1],\cdots,\mathbf{y}[N-1]$的线性组合，系数的求解可以通过多项式快速幂算法来实现，复杂度为$O(N^2 \log k)$；矩阵乘法的复杂度为$O(N^2)$。因此每求一个$\mathbf{y}[k]$的复杂度为$O(N^2 \log k)$。

对于非齐次方程，由于需要计算$\sum_{i=0}^{k-1} A^{k-1-i} \mathbf{f}[i]$，如果每个$A^{k-1-i}$都通过多项式快速幂算法来计算，那么每求一个$\mathbf{y}[k]$的复杂度为$O(k N^2 \log k)$，此时复杂度比直接迭代还要大。

将上述三种算法在齐次/非齐次方程下的复杂度汇总如下：
\begin{table}[h]
    \centering
    \begin{tabular}{|c|c|c|}
        \hline
        算法 & 齐次方程 & 非齐次方程 \\
        \hline
        直接迭代 & $O(k N^2)$ & $O(k N^2)$ \\
        \hline
        Jordan分解算法 & $O(N^3)$ & $O(N^3)+O(k N^2)$ \\
        \hline
        Cayley-Hamilton定理递推算法 & $O(N^3)+O(N^2 \log k)$ & $O(N^3)+O(k N^2 \log k)$ \\
        \hline
    \end{tabular}
\end{table}

可见，当只需求解较小的$k$或方程含有非齐次项时，直接迭代算法的效率更高；当$k\gg N$且方程齐次时，Jordan分解算法效率最高。但是，Jordan分解算法的复杂度中$N^3$的系数较大，有时使用Cayley-Hamilton定理递推算法会更快一些。

\subsection{解矩阵}\label{sub:fundamental_solution_matrix}

类比微分方程组，定义：
\begin{definition}[解矩阵]
    常系数线性差分方程组
    $$\mathbf{y}[n+1] = A \mathbf{y}[n] + \mathbf{f}[n]$$
    的解矩阵为
    $$\phi[n]=A^n$$
\end{definition}

得到这个定义的思路有两种。一种思路是考虑将齐次差分方程中的列向量$\mathbf{y}[n]$替换为一个矩阵值函数，则$\phi[n]=A^n$恰好满足方程\begin{align*}
    \phi[n+1] = A \phi[n],\phi[0] = I
\end{align*}

另一种思路是类比微分方程。常系数线性微分方程最典型的解的形式为$e^{\lambda t}$，而方程组的解矩阵为$\phi(t)=e^{At}$；对于差分方程，将$\lambda^n$中的$\lambda$替换为矩阵$A$，就得到解矩阵$\phi[n]=A^n$。

解矩阵的一个重要特征是，\textbf{它的列向量是齐次差分方程的解}，因此方程的齐次解总能表示为$\phi[n]\mathbf{c}$的形式，其中$\mathbf{c}\in\mathbb{R}^N$是一个常向量。与微分方程不同的是，当$A$不可逆时，$\phi[n]$也是不可逆的，因此$\phi^{-1}[n]$并不总是存在的。

当$A$可逆时，有$\phi^{-1}[n] = A^{-n}=\phi[-n]$，它的列空间有$N$维，恰为差分方程组的解空间，此时称之为\textbf{基本解矩阵}。

事实上，对于变系数的线性差分方程组\begin{align*}
    \mathbf{y}[n+1] = A[n] \mathbf{y}[n]+ \mathbf{f}[n]
\end{align*}
我们同样可以定义解矩阵：
\begin{align*}
    \phi[n] = A[n-1] A[n-2] \cdots A[0]
\end{align*}
它满足$\phi[n+1] = A[n] \phi[n],\phi[0] = I$，因此其列向量确为方程组的解。但由于$A[n]$不为常数矩阵，我们通常是求不出$\phi[n]$的解析表达式的，这一点与\ref{sec:Equation_System}中的线性微分方程组有些类似。

解矩阵在差分方程组理论中，仍然扮演着“基本解”的角色。以非齐次项$\mathbf{f}[n]$为激励，方程组\begin{align*}
    \mathbf{y}[n+1] = A \mathbf{y}[n] + \mathbf{f}[n],\mathbf{y}_0 = \mathbf{y}[0]
\end{align*}
的解$\mathbf{y}[n]$为响应，则方程描述了一个离散的因果线性时不变系统，类比\ref{sub:constant_coefficient_linear_equation_system}中的常系数线性微分方程组，可以猜想零状态响应可以表示为卷积和，然而，我们实际上应该使用$\phi[n-1]u[n-1]$做卷积才能得到正确的结论，我们暂时不做解释。\begin{align*}
    \mathbf{y}_{zs}[n] = (\phi[n-1]u[n-1]) * \mathbf{f}[n]=\sum_{i=0}^{n-1} \phi[n-1-i]\mathbf{f}[i]=\sum_{i=0}^{n-1} A^{n-1-i} \mathbf{f}[i]
\end{align*}
求和范围的改变是因为我们认为激励$\mathbf{f}[n]$是在$n=0$时刻施加于系统的。这与上一小节中用迭代法得到的通解公式一致。

零输入响应是方程的齐次解。由于解矩阵的列向量是差分方程的解，零输入响应可以表示为解矩阵的列向量的线性组合：\begin{align*}
    \mathbf{y}_{zi}[n] = \phi[n]\mathbf{c},\mathbf{c}\in\mathbb{R}^N\tag{$*$}
\end{align*}

代入初始值$\mathbf{y}[0]=\mathbf{y}_0$，可以求出$\mathbf{c}$的表达式:
\begin{align*}
    \mathbf{y}_0 = \phi[0]\mathbf{c}\implies \mathbf{c} = \phi^{-1}[0]\mathbf{y}_0 = \mathbf{y}_0
\end{align*}
尽管$\phi[n]$不一定可逆，但在$n=0$时，$\phi[0]=I$。代入$(*)$式，就得到零输入响应的表达式：
\begin{align*}
    \mathbf{y}_{zi}[n] = \phi[n]\mathbf{c} = \phi[n]\mathbf{y}_0
\end{align*}
最终，我们得到了常系数线性差分方程组的通解公式，它与\ref{sub:constant_coefficient_linear_equation_system}中的常系数线性微分方程组的通解公式非常相似，并且与上一小节中通过迭代法得到的通解公式一致：
\begin{theorem}
    设$\phi[n]$为带有初始值的方程组\begin{align*}
        \mathbf{y}[n+1] = A \mathbf{y}[n]+\mathbf{f}[n],\mathbf{y}_0 = \mathbf{y}[0]
    \end{align*}
    的解矩阵，则方程组的通解为\begin{align*}
        \mathbf{y}[n]& = \phi[n]\mathbf{y}_0 + \sum_{i=0}^{n} \phi[i]\mathbf{f}[n-i]\\
        & = A^n \mathbf{y}_0 + \sum_{i=0}^{n-1} A^{n-1-i} \mathbf{f}[i]
    \end{align*}
\end{theorem}

\subsection{常系数线性差分方程组的单边z变换解法}\label{sub:difference_equation_system_z}

单边z变换同样能够处理常系数线性差分方程组，并且能够一并处理初始值。我们首先给出离散矩阵值函数的单边z变换的定义：

\begin{definition}[离散矩阵值函数的单边z变换]
    设$A[n]$为$N\times N$的矩阵值函数，则\textbf{单边z变换}就是它的各个分量的单边z变换：
    \begin{align*}
        A(z) = \sum_{n=0}^{\infty} A[n] z^{-n}
    \end{align*}
    收敛域为各个分量的单边z变换的收敛域的交集。
\end{definition}

特别地，列向量可以视作只有一列的矩阵，因此向量值函数的单边z变换也是良定义的。

考虑如下常系数线性差分方程组\begin{align*}
    \mathbf{y}[n+1] = A \mathbf{y}[n] + \mathbf{f}[n],\mathbf{y}_0 = \mathbf{y}[0]
\end{align*}
两边同时乘以$u[n]$并进行单边z变换，得到\begin{align*}
    z\mathbf{Y}(z) - z\mathbf{y}_0 = A \mathbf{Y}(z) + \mathbf{F}(z)
\end{align*}
其中$\mathbf{Y}(z)$和$\mathbf{F}(z)$分别是$\mathbf{y}[n]$和$\mathbf{f}[n]$的单边z变换。因此，\begin{align*}
    (zI - A) \mathbf{Y}(z) = \mathbf{F}(z) + z\mathbf{y}_0
\end{align*}
两边同时求逆z变换，即可得到方程的解：\begin{align*}
    \mathbf{y}[n] = \mathcal{Z}^{-1} \left\{(zI - A)^{-1} (\mathbf{F}(z) + z\mathbf{y}_0)\right\}\tag{$**$}
\end{align*}

令$\Phi(z)=(zI - A)^{-1}$，称之为差分方程组的\textbf{特征矩阵}或\textbf{预解矩阵}。注意到\begin{align*}
    \Phi(z)=(zI - A)^{-1} = z^{-1} (I - z^{-1} A)^{-1} = z^{-1} \sum_{i=0}^{\infty}A^i z^{-i} = z^{-1}\mathcal{Z}\{\phi[n]u[n]\}
\end{align*}
因此$z\Phi(z)$就是解矩阵$\phi[n]$的单边z变换（注意，这与微分方程组不同）。代入$(**)$式，得到
\begin{align*}
    \mathbf{y}[n] & = \mathcal{Z}^{-1} \left\{\Phi(z)(\mathbf{F}(z) + z\mathbf{y}_0)\right\}\\
    &= \mathcal{Z}^{-1} \left\{\Phi(z) \mathbf{F}(z)\right\} + \mathcal{Z}^{-1} \left\{z\Phi(z)\mathbf{y}_0\right\}\\
    & = \mathcal{Z}^{-1} \left\{z^{-1}\mathcal{Z}\{\phi[n]u[n]\} \mathbf{F}(z)\right\} +\phi[n]u[n]\mathbf{y}_0
\end{align*}
利用单边z变换的移位性质以及卷积定理，得到\begin{align*}
    \mathcal{Z}^{-1} \left\{z^{-1}\mathcal{Z}\{\phi[n]u[n]\} \mathbf{F}(z)\right\}&=(\phi[n-1]u[n-1]) * (\mathbf{f}[n]u[n])\\
    &=\sum_{i=0}^{n} \phi[n-1-i] \mathbf{f}[i]\\
    & = \sum_{i=0}^{n} A^{n-1-i} \mathbf{f}[i]
\end{align*}

综上，我们得到了常系数线性差分方程组的通解公式，它与之前推导的公式相一致：
\begin{align*}
    \mathbf{y}[n]& = \phi[n]\mathbf{y}_0 + \sum_{i=0}^{n} \phi[n-1-i]\mathbf{f}[i]\\
    & = A^n \mathbf{y}_0 + \sum_{i=0}^{n} A^{n-1-i} \mathbf{f}[i]
\end{align*}

\begin{example}
    求解常系数线性差分方程组在$n\in\mathbb{N}$上的解：
    \begin{align*}
        \begin{bmatrix}
            y_1[n+1]\\
            y_2[n+1]
        \end{bmatrix}=\begin{bmatrix}
            1 & 1\\
            0 & 1
        \end{bmatrix}
        \begin{bmatrix}
            y_1[n]\\
            y_2[n]
        \end{bmatrix} + \begin{bmatrix}
            0\\
            u[n]
        \end{bmatrix},\begin{bmatrix}
            y_1[0]\\
            y_2[0]
        \end{bmatrix} = \begin{bmatrix}
            1\\
            0
        \end{bmatrix}
    \end{align*}

    解：记
    $$A=\begin{bmatrix}
            1 & 1\\
            0 & 1
        \end{bmatrix}$$
    两边同时乘以$u[n]$并进行单边z变换，得到\begin{align*}
        z\begin{bmatrix}
            Y_1(z) \\
            Y_2(z)
        \end{bmatrix} - z\begin{bmatrix}
            1\\
            0
        \end{bmatrix} = A\begin{bmatrix}
            Y_1(z)\\
            Y_2(z)
        \end{bmatrix} + \begin{bmatrix}
            0\\
            \dfrac{z}{z-1}
        \end{bmatrix}
    \end{align*}
    从而\begin{align*}
        (zI - A) \begin{bmatrix}
            Y_1(z)\\
            Y_2(z)
        \end{bmatrix} = \begin{bmatrix}
            z\\
            \dfrac{z}{z-1}
        \end{bmatrix}
    \end{align*}
    计算特征矩阵：\begin{align*}
        \Phi(z)=(zI - A)^{-1} =\begin{bmatrix}
            z-1 & -1\\
            0 & z-1
        \end{bmatrix}^{-1}= \begin{bmatrix}
            \dfrac{1}{z-1} & \dfrac{1}{(z-1)^2}\\
            0 & \dfrac{1}{z-1}
        \end{bmatrix}
    \end{align*}
    因此\begin{align*}
        \begin{bmatrix}
            Y_1(z)\\
            Y_2(z)
        \end{bmatrix} & = \Phi(z) \begin{bmatrix}
            z\\
            \dfrac{z}{z-1}
        \end{bmatrix}\\
        & = \begin{bmatrix}
            \dfrac{1}{z-1} & \dfrac{1}{(z-1)^2}\\
            0 & \dfrac{1}{z-1}
        \end{bmatrix} \begin{bmatrix}
            z\\
            \dfrac{z}{z-1}
        \end{bmatrix}\\
        & = \begin{bmatrix}
            \dfrac{z}{z-1} + \dfrac{z}{(z-1)^3}\\
            \dfrac{z}{(z-1)^2}
        \end{bmatrix}
    \end{align*}
    最后，选取半径大于$1$的围道$C$，求逆z变换：\begin{align*}
        y_1[n] & = \dfrac{1}{2\pi i} \oint_C \left(\dfrac{z}{z-1} + \dfrac{z}{(z-1)^3}\right) z^{n-1} dz\\
        & = \mathrm{Res}\left(\dfrac{z^n}{z-1},1\right) + \mathrm{Res}\left(\dfrac{z^n}{(z-1)^3},1\right)\\
        &=\lim_{z\to 1}z^n+ \lim_{z\to 1} \dfrac{1}{2}\dfrac{d^2}{dz^2}z^n\\
        & = 1 + \dfrac{n(n+1)}{2}\\
        y_2[n] & = \dfrac{1}{2\pi i} \oint_C \dfrac{z}{(z-1)^2} z^{n-1} dz\\
        & = \mathrm{Res}\left(\dfrac{z^n}{(z-1)^2},1\right)\\
        &=\lim_{z\to 1} \dfrac{d}{dz} z^n\\
        & = n
    \end{align*}
\end{example}

\subsection{*离散系统的状态变量分析}\label{sub:discrete_state_variable}


